% Generated draft: requires mathematical and source review.
\documentclass[11pt]{article}
\usepackage[margin=25mm]{geometry}
\usepackage{fontspec}
\setmainfont{Latin Modern Roman}
\usepackage{amsmath,amssymb,mathtools,unicode-math}
\setmathfont{Latin Modern Math}
\usepackage{enumitem,needspace,xurl,hyperref}
\setlength{\parindent}{0pt}
\setlength{\parskip}{5pt}
\setlength{\emergencystretch}{3em}
\newcommand{\N}{\mathbb{N}}
\newcommand{\Z}{\mathbb{Z}}
\newcommand{\Q}{\mathbb{Q}}
\newcommand{\R}{\mathbb{R}}
\newcommand{\C}{\mathbb{C}}
\newcommand{\F}{\mathbb{F}}
\newcommand{\A}{\mathbb{A}}
\newcommand{\PP}{\mathbb{P}}
\newcommand{\E}{\mathbb{E}}
\newcommand{\eps}{\varepsilon}
\newcommand{\dd}{\,\mathrm{d}}
\begin{document}
\section*{UnsolvedMath ID: 20000393}
\textbf{Source problem code:} catalogue entry
\subsection{Definitions and mathematical statement}
Jordan Ellenberg asks which algebro-geometric properties of elliptic surfaces explain the Poonen--Rains and BKLPR distribution models. In a function-field formulation, an elliptic surface over a curve over $\F_q$ has generic fiber an elliptic curve over that curve's function field. A proposed explanation should specify the family, height ordering, Selmer level, geometric realization, monodromy and limiting distribution. None of these choices or the precise model law is given in the original.

The attached unreviewed aid reports substantial function-field progress: finite \'etale Selmer spaces, a quadratic intersection form on primitive elliptic-surface cohomology, large orthogonal monodromy and Frobenius equidistribution produce random orthogonal fixed-space kernels. It attributes a stable identification with the finite-level BKLPR law to Feng--Landesman--Rains, in a limit taking large $q$ first and then height. The field/characteristic restrictions, family and exact two-stage limit are not reproduced; this does not assert an arbitrary joint limit or a number-field theorem.

Its separate linear-algebra bridge fixes an odd prime $\ell$ and a nondegenerate finite-dimensional quadratic space $(V,Q)$ over $\F_\ell$. For every $g\in O(V,Q)$, set
\[
\Delta=\{(v,v):v\in V\},\qquad\Gamma_g=\{(v,gv):v\in V\}.
\]
The aid claims that these are maximal isotropic subspaces of $(V\oplus V,Q\oplus(-Q))$, that $\Delta\cap\Gamma_g\simeq\ker(g-1)$, and that
\[
(-1)^{\dim\ker(g-1)}=(-1)^{\dim V}\det(g).
\]
The determinant is identified with its sign $\pm1$; no semisimplicity assumption on $g$ is imposed. If $V$ contains a hyperbolic plane, the aid additionally claims
\[
\frac1{|O(V,Q)|}\sum_{g\in O(V,Q)}|\ker(g-1)|=\ell+1.
\]
A hyperbolic plane is a two-dimensional nondegenerate quadratic subspace with isotropic lines. Uniform sampling here is from the full orthogonal group, including both determinant components; applying the mean to an arithmetic subgroup or Frobenius coset needs a separate comparison. These finite-space statements illustrate an intersection-of-maximal-isotropics mechanism, but do not themselves prove the original geometric distribution assertion.
\subsection{Short English statement}
Explain random Selmer models through elliptic-surface geometry, cohomological quadratic forms, monodromy and Frobenius statistics. The aid supplies a finite-orthogonal-space bridge while retaining the specific function-field limit and group-sampling qualifications.
\subsection{Sources}
\begin{itemize}
\item Jordan Ellenberg, Question 15, AIM \emph{Arithmetic statistics over finite fields and function fields}: \url{https://aimath.org/pastworkshops/arithstatffieldproblems.pdf}. Canonical location: \texttt{aim-algebraic-number-theory-notes.json}, notes[44]. The original is labeled exact; remarks and literature are empty. No PDF was consulted.
\item Original model references: Poonen--Rains and BKLPR, without bibliographic entries or a precise probability law. The aid names Feng--Landesman--Rains for the function-field identification, also without a complete reference.
\item The attached graph-diagonal aid is machine generated and not peer reviewed, with low novelty confidence. Its published-work assessment and candidate lemmas are not independently verified here.
\item Local registry: \texttt{data/unsolvedmath/problems.json}, permanent catalogue entry, exact source code \texttt{catalogue entry}.
\end{itemize}
\end{document}
